El repositorio \texttt{FeedTheRealm-org/Documents} centraliza la documentación académica y técnica del proyecto Feed the Realm
en formato \LaTeX. Su principal característica es la compilación y publicación automática de los documentos generados: 
cada vez que se etiqueta una nueva versión en el repositorio, un \textit{pipeline} de integración continua compila los 
fuentes \LaTeX{} y publica los PDFs resultantes como \textit{GitHub Releases}.
\vspace{0.25cm}

El repositorio está organizado de la siguiente manera:
\vspace{0.25cm}

\begin{itemize}
    \item \textbf{Carpetas de documentos} (p.\,ej.\ \texttt{proposal/}, \texttt{simulate/}): cada una contiene los fuentes 
    \LaTeX{} de un documento, con \texttt{main.tex} como punto de entrada.
    \item \textbf{\texttt{files.yaml}}: archivo de registro donde se declaran todos los documentos a compilar, especificando 
    nombre, descripción y directorio de entrada.
    \item \textbf{\texttt{Makefile}}: expone los comandos principales para construir el entorno y compilar los documentos.
    \item \textbf{\texttt{Dockerfile}}: define la imagen Docker que contiene el entorno \LaTeX{} necesario para la compilación, 
    evitando dependencias locales.
    \item \textbf{\texttt{generar\_docs.py}}: script Python auxiliar para la generación de documentos.
    \item \textbf{\texttt{.github/workflows/}}: contiene las definiciones del \textit{pipeline} de CI/CD que automatiza 
    la compilación y publicación.
\end{itemize}

Para agregar un nuevo documento al repositorio se deben seguir los siguientes pasos:
\vspace{0.25cm}

\begin{enumerate}
    \item Crear una carpeta con los fuentes \LaTeX{} del documento, utilizando \texttt{main.tex} como punto de entrada.
    \item Registrar el documento en \texttt{files.yaml} con su nombre, descripción y ruta al directorio correspondiente.
    \item Crear y publicar un \textit{tag} de versión para disparar el \textit{pipeline} de compilación y publicación:
\end{enumerate}

\begin{verbatim}
git tag -a v1.0.0 -m "Release version 1.0.0"
git push --tags
\end{verbatim}

Una vez ejecutado el \textit{pipeline}, los PDFs compilados quedan disponibles públicamente en la sección de \textit{Releases} 
del repositorio, accesibles mediante una URL de la forma:

\begin{verbatim}
https://github.com/FeedTheRealm-org/Documents/releases/tag/v1.0.0
\end{verbatim}

La compilación se realiza dentro de un contenedor Docker para no requerir una instalación local de \LaTeX{}. Los comandos 
disponibles a través del \texttt{Makefile} son:
\vspace{0.25cm}

\begin{verbatim}
make image       # Construye la imagen Docker con el entorno LaTeX
make run         # Compila todos los documentos definidos en files.yaml
make run-<name>  # Compila únicamente el documento especificado
make clean       # Elimina todos los PDFs generados del directorio raíz
\end{verbatim}

Este enfoque garantiza reproducibilidad en cualquier entorno de desarrollo, independientemente del sistema operativo o la 
distribución de \LaTeX{} instalada localmente.
\vspace{0.25cm}
